% CP-VI-0002
% target_chapter: VI/07 — Spectrum of a Ring
% source_unit_id: IMP-DL-0E3E7EC766-P002
% source: Downloads/theory-of-algebraic-geometry-1.tex L358-L711
% solution_provenance: SOURCE_EXPLICIT_SOLUTION

\subsection*{Problem 2}
	
	Describe the topological spaces
	\[
	\operatorname{Spec} \mathbb R[x],
	\qquad
	\operatorname{Spec} \mathbb C[x,y],
	\qquad
	\operatorname{Spec} \mathbb Z[x],
	\qquad
	\operatorname{Spec} \mathbb F_p[x],
	\]
	where \(\mathbb F_p\) is the field with \(p\) elements.
	
	\bigskip
	
	\textbf{Solution.}
	
	We describe the points and the closure relations in each case.
	
	\subsection*{1. The space \(\operatorname{Spec} \mathbb R[x]\)}
	
	The ring \(\mathbb R[x]\) is a principal ideal domain. Therefore its prime ideals
	are
	\[
	(0)
	\]
	and the ideals generated by irreducible polynomials over \(\mathbb R\).
	
	Over \(\mathbb R\), the irreducible polynomials are of two types:
	
	\begin{enumerate}[label=\arabic*)]
		\item linear polynomials,
		\[
		x-a,
		\qquad a\in \mathbb R;
		\]
		
		\item quadratic polynomials with negative discriminant.
	\end{enumerate}
	
	Hence the points of \(\operatorname{Spec}\mathbb R[x]\) are
	\[
	(0),
	\]
	the ideals
	\[
	(x-a),
	\qquad a\in \mathbb R,
	\]
	and the ideals
	\[
	(q(x)),
	\]
	where \(q(x)\in\mathbb R[x]\) is an irreducible quadratic polynomial.
	
	The point \((0)\) is the generic point of the whole space. Indeed,
	\[
	\overline{\{(0)\}}=\operatorname{Spec}\mathbb R[x].
	\]
	This happens because
	\[
	V((0))=\operatorname{Spec}\mathbb R[x].
	\]
	
	Every nonzero prime ideal in \(\mathbb R[x]\) is maximal, because \(\mathbb R[x]\)
	is a PID and any nonzero prime ideal is generated by an irreducible polynomial.
	Thus all nonzero prime ideals are closed points.
	
	Therefore the closed points are
	\[
	(x-a),
	\qquad a\in\mathbb R,
	\]
	and
	\[
	(q(x)),
	\]
	where \(q(x)\) is irreducible quadratic over \(\mathbb R\).
	
	Geometrically, \(\operatorname{Spec}\mathbb R[x]\) is the affine line over \(\mathbb R\), but
	with closed points corresponding not only to real roots, but also to conjugate
	pairs of complex non-real roots.
	
	The points \((x-a)\) correspond to real points
	\[
	a\in\mathbb R.
	\]
	The quadratic ideals correspond to conjugate pairs
	\[
	\alpha,\overline{\alpha}\in\mathbb C\setminus \mathbb R.
	\]
	
	Thus \(\operatorname{Spec}\mathbb R[x]\) is a one-dimensional topological space with one
	generic point and many closed points.
	
	The proper closed subsets are finite sets of closed points.
	
	\subsection*{2. The space \(\operatorname{Spec} \mathbb C[x,y]\)}
	
	The ring \(\mathbb C[x,y]\) is the coordinate ring of the affine plane over
	\(\mathbb C\). Thus \(\operatorname{Spec}\mathbb C[x,y]\) is the affine plane
	\[
	\mathbb A^2_{\mathbb C}
	\]
	together with generic points for irreducible algebraic subsets.
	
	The ring \(\mathbb C[x,y]\) has Krull dimension \(2\). Therefore its prime ideals
	have height \(0\), \(1\), or \(2\).
	
	There is one height \(0\) prime ideal:
	\[
	(0).
	\]
	This is the generic point of the whole affine plane. Its closure is
	\[
	\overline{\{(0)\}}=\operatorname{Spec}\mathbb C[x,y].
	\]
	
	The height \(1\) prime ideals are of the form
	\[
	(f),
	\]
	where \(f\in\mathbb C[x,y]\) is an irreducible polynomial. These points are not
	closed. Each such point is the generic point of the irreducible algebraic curve
	defined by
	\[
	f(x,y)=0.
	\]
	Its closure is
	\[
	\overline{\{(f)\}}=V(f).
	\]
	
	The height \(2\) prime ideals are maximal ideals. By Hilbert's Nullstellensatz,
	the maximal ideals of \(\mathbb C[x,y]\) are exactly
	\[
	(x-a,y-b),
	\qquad (a,b)\in\mathbb C^2.
	\]
	These are the closed points.
	
	Therefore the points of \(\operatorname{Spec}\mathbb C[x,y]\) are:
	
	\begin{enumerate}[label=\arabic*)]
		\item the generic point
		\[
		(0);
		\]
		
		\item generic points of irreducible curves,
		\[
		(f),
		\qquad f\in\mathbb C[x,y]\text{ irreducible};
		\]
		
		\item closed points,
		\[
		(x-a,y-b),
		\qquad (a,b)\in\mathbb C^2.
		\]
	\end{enumerate}
	
	The closure relations are:
	\[
	\overline{\{(0)\}}=\operatorname{Spec}\mathbb C[x,y],
	\]
	\[
	\overline{\{(f)\}}=V(f),
	\]
	and
	\[
	\overline{\{(x-a,y-b)\}}=\{(x-a,y-b)\}.
	\]
	
	Thus \(\operatorname{Spec}\mathbb C[x,y]\) is the affine plane, but topologically it contains
	more than ordinary complex points. It also contains generic points for curves
	and a generic point for the whole plane.
	
	\subsection*{3. The space \(\operatorname{Spec} \mathbb Z[x]\)}
	
	The ring \(\mathbb Z[x]\) has dimension \(2\). It may be viewed as the coordinate
	ring of the affine line over \(\operatorname{Spec}\mathbb Z\). Therefore
	\[
	\operatorname{Spec}\mathbb Z[x]
	\]
	is an arithmetic surface.
	
	There is one height \(0\) prime:
	\[
	(0).
	\]
	This is the generic point of the whole space:
	\[
	\overline{\{(0)\}}=\operatorname{Spec}\mathbb Z[x].
	\]
	
	The height \(1\) primes include two important types.
	
	First, for every prime number \(p\), we have the prime ideal
	\[
	(p)\subseteq \mathbb Z[x].
	\]
	Indeed,
	\[
	\mathbb Z[x]/(p)\cong \mathbb F_p[x],
	\]
	which is an integral domain. Therefore \((p)\) is prime.
	
	The closed subset
	\[
	V(p)\subseteq \operatorname{Spec}\mathbb Z[x]
	\]
	is the fiber over the prime \(p\). Since
	\[
	\mathbb Z[x]/(p)\cong \mathbb F_p[x],
	\]
	we have
	\[
	V(p)\cong \operatorname{Spec} \mathbb F_p[x].
	\]
	These are called vertical curves.
	
	Second, if \(F(x)\in \mathbb Z[x]\) is primitive and irreducible, then
	\[
	(F(x))
	\]
	is a height \(1\) prime ideal. These are called horizontal curves, because they
	spread over \(\operatorname{Spec}\mathbb Z\).
	
	The maximal ideals of \(\mathbb Z[x]\) are the closed points. They are of the
	form
	\[
	(p,\overline f(x)),
	\]
	where \(p\) is a prime number and \(\overline f(x)\in\mathbb F_p[x]\) is an
	irreducible polynomial.
	
	Equivalently, if \(f(x)\in\mathbb Z[x]\) is a lift of \(\overline f(x)\), then
	the maximal ideal may be written as
	\[
	(p,f(x)).
	\]
	The residue field is
	\[
	\mathbb Z[x]/(p,f(x))
	\cong
	\mathbb F_p[x]/(\overline f(x)),
	\]
	which is a finite field.
	
	Therefore the points of \(\operatorname{Spec}\mathbb Z[x]\) are:
	
	\begin{enumerate}[label=\arabic*)]
		\item the generic point
		\[
		(0);
		\]
		
		\item vertical curve points
		\[
		(p),
		\]
		where \(p\) is a prime number;
		
		\item horizontal curve points
		\[
		(F(x)),
		\]
		where \(F(x)\in\mathbb Z[x]\) is primitive and irreducible;
		
		\item closed points
		\[
		(p,\overline f(x)),
		\]
		where \(p\) is prime and \(\overline f(x)\in\mathbb F_p[x]\) is irreducible.
	\end{enumerate}
	
	The closure relations are as follows.
	
	The closure of the generic point is the whole space:
	\[
	\overline{\{(0)\}}=\operatorname{Spec}\mathbb Z[x].
	\]
	
	The closure of a height \(1\) point is an irreducible curve:
	\[
	\overline{\{(p)\}}=V(p)\cong \operatorname{Spec}\mathbb F_p[x],
	\]
	and
	\[
	\overline{\{(F)\}}=V(F).
	\]
	
	The maximal ideals are closed points:
	\[
	\overline{\{(p,\overline f)\}}=\{(p,\overline f)\}.
	\]
	
	Thus \(\operatorname{Spec}\mathbb Z[x]\) is a two-dimensional arithmetic space. It has a
	generic point, one-dimensional curve points, and closed points with finite
	residue fields.
	
	\subsection*{4. The space \(\operatorname{Spec} \mathbb F_p[x]\)}
	
	The ring \(\mathbb F_p[x]\) is a principal ideal domain. Therefore its prime
	ideals are
	\[
	(0)
	\]
	and the ideals generated by irreducible polynomials.
	
	Hence the points of \(\operatorname{Spec}\mathbb F_p[x]\) are
	\[
	(0)
	\]
	and
	\[
	(f(x)),
	\]
	where \(f(x)\in\mathbb F_p[x]\) is irreducible.
	
	The point \((0)\) is the generic point:
	\[
	\overline{\{(0)\}}=\operatorname{Spec}\mathbb F_p[x].
	\]
	
	Every nonzero prime ideal in \(\mathbb F_p[x]\) is maximal. Therefore each
	\[
	(f(x))
	\]
	with \(f\) irreducible is a closed point.
	
	The linear ideals
	\[
	(x-a),
	\qquad a\in \mathbb F_p,
	\]
	correspond to the \(p\) rational points of the affine line over \(\mathbb F_p\).
	
	If \(f(x)\) is irreducible of degree \(d>1\), then the residue field is
	\[
	\mathbb F_p[x]/(f(x))\cong \mathbb F_{p^d}.
	\]
	Thus irreducible polynomials of higher degree correspond to closed points whose
	residue fields are finite extensions of \(\mathbb F_p\).
	
	Therefore \(\operatorname{Spec}\mathbb F_p[x]\) is the affine line over \(\mathbb F_p\), but
	its closed points are not only the \(p\) rational points. They also include
	points corresponding to irreducible polynomials of all degrees.
	
	As in the case of a polynomial ring in one variable over a field, the proper
	closed subsets are finite sets of closed points.
